\documentclass[12pt]{article}
\usepackage[T1]{fontenc}
\usepackage[utf8]{inputenc}
\usepackage{lmodern}
\usepackage{textcomp}
\usepackage{enumitem}
\usepackage{geometry}
\geometry{margin=1in}
\begin{document}
\begin{center}
Answer Key - Philosophy - Graduate Level Assessment 19 \
\small (Answers are ordered exactly as in the question paper.)
\end{center}

\section*{Multiple Choice}
\begin{enumerate}[label=\arabic*.]
\item \textbf{Answer:} A
\\ \textit{Options:} A) preference.; B) natural law.; C) self-interest.; D) maximizing utility.; E) ethical relativism.
\\ \textit{Note:} Anscombe argues that moral obligation is rooted in human preferences and desires, suggesting that our moral judgments are closely linked to what individuals value and choose to prioritize in their lives.
\item \textbf{Answer:} C
\\ \textit{Options:} A) elements of nature that do not exist independently; B) only things existing apart from our minds; C) only sensations existing in our minds; D) physical objects; E) manifestations of our subconscious
\\ \textit{Note:} Berkeley's philosophy posits that heat and cold, like all qualities of objects, are not inherent properties of the external world but rather subjective sensations experienced in the mind, highlighting his idealist view that reality consists of ideas rather than material substances.
\item \textbf{Answer:} A
\\ \textit{Options:} A) not essential to our existence; B) frequently used in daily life; C) not visible to the human eye; D) only appreciated by experts; E) universally liked
\\ \textit{Note:} The correct answer is based on the understanding that aesthetics focuses on the appreciation and evaluation of beauty and art, which are considered non-essential aspects of human experience, distinct from the necessities required for survival.
\item \textbf{Answer:} A
\\ \textit{Options:} A) dynamic use.; B) descriptive use.; C) inferential use.; D) propositional use.; E) expressive use.
\\ \textit{Note:} The word "good" possesses a dynamic quality in its emotive meaning, allowing it to adapt and resonate in various contexts, which aligns with Stevenson's emphasis on the importance of emotive language in effective communication.
\item \textbf{Answer:} D
\\ \textit{Options:} A) a list of crimes ranked according to their seriousness.; B) a scale of punishments that correspond to the seriousness of certain crimes.; C) treating criminals humanely.; D) the death penalty for the most serious crimes.
\\ \textit{Note:} Nathanson's version of retributivism advocates for proportional punishment based on moral desert rather than necessitating the death penalty, emphasizing that not all serious crimes merit the most extreme form of punishment.
\end{enumerate}

\section*{True / False}
\begin{enumerate}[label=\arabic*.]
\setcounter{enumi}{5}
\item \textbf{Answer:} False
\\ \textit{Options:} A) True; B) False
\\ \textit{Note:} Dershowitz argues that while psychological manipulation techniques can be ethically problematic, there are contexts, such as in defense of civil liberties or in certain therapeutic settings, where they may be considered acceptable or necessary.
\item \textbf{Answer:} False
\\ \textit{Options:} A) True; B) False
\\ \textit{Note:} Anscombe argues that moral obligation is grounded in ethical principles and human reason, rather than being solely determined by the legal codes of a country.
\item \textbf{Answer:} True
\\ \textit{Options:} A) True; B) False
\\ \textit{Note:} Hasty generalization is considered a logical fallacy because it occurs when a conclusion is reached based on a small or unrepresentative sample, failing to provide adequate evidence to support the broader claim.
\item \textbf{Answer:} True
\\ \textit{Options:} A) True; B) False
\\ \textit{Note:} In Japanese culture, particularly within Shinto and certain Buddhist traditions, the belief in divine benevolence and the rewards of devotion is a central theme that emphasizes the importance of harmonious relationships between humans and the divine.
\item \textbf{Answer:} False
\\ \textit{Options:} A) True; B) False
\\ \textit{Note:} Taurek argues that the moral value of saving one individual can outweigh the aggregate value of saving multiple individuals, emphasizing the significance of individual lives and personal connections rather than a strict utilitarian calculus.
\end{enumerate}

\section*{Long Form Response}
\begin{enumerate}[label=\arabic*.]
\setcounter{enumi}{10}
\item \textbf{Answer:} Michael Walzer identifies four main justifications for terrorism: as a response to oppression, as a means of achieving political goals, as a tactic of war, and as a last resort when other avenues are unavailable. However, he argues that terrorism fundamentally undermines the social contract and moral framework necessary for political discourse. The notion of terrorism as a form of freedom of speech is problematic within Walzer's framework because it fails to recognize the inherent violence and dehumanization involved in terrorist acts, which obliterate the possibility of rational dialogue and mutual respect. Thus, while expressing dissent is a critical aspect of free speech, the violent imposition of one's will through terrorism contradicts the very principles of communication and understanding that free speech encompasses.
\\ \textit{Note:} correct because it accurately identifies Walzer's justifications for terrorism while highlighting how the concept of terrorism as a form of freedom of speech contradicts his belief in the necessity of a moral framework for political discourse, which terrorism undermines through its inherent violence.
\item \textbf{Answer:} A slippery-slope argument is considered fallacious when it lacks sufficient evidence to demonstrate that a particular action will inevitably lead to a chain of undesirable consequences. For instance, arguing that legalizing marijuana will inevitably result in the legalization of all drugs fails if it does not adequately establish a causal connection between these outcomes. The fallacy arises from the assumption that the initial action will unavoidably lead to extreme results without empirical support. Thus, while some slippery-slope arguments may have merit when plausible connections exist, those that rely solely on fear of consequences without substantiation are logically flawed and unconvincing.
\\ \textit{Note:} correct because it identifies that a slippery-slope argument is fallacious when it lacks sufficient evidence for a causal connection between the initial action and its extreme consequences, as illustrated by the example of legalizing marijuana leading to the legalization of all drugs without empirical support.
\end{enumerate}

\vfill
\noindent\textit{End of Answer Key}
\end{document}